%%%%

%%%   Самарский государственный технический университет

%%%   Кафедра прикладной математики и информатики

%%%

%%%   Преамбула для статьи в журнал "Вестник Самарского государственного

%%%   технического университета. Серия: «Физико-математические науки»"

%%%   Ориентирована под MikTEx 2.4 for Win32

%%%

%%%   Хотя сам журнал собирается под GNU/Linux на дистрибутиве TeXLive



\documentclass[11pt,a4paper,twoside]{article}



\usepackage[cp1251]{inputenc}

\usepackage[T2A]{fontenc}

\usepackage[english,russian]{babel}

\usepackage{samgtubib}

\usepackage[dvips]{graphicx}

\usepackage[multiple]{footmisc}



\usepackage{samgtu}

\renewcommand{\VestnikYear}{2009} % Год издания Вестника

\renewcommand{\VestnikNumber}{2\,(19)} % Номер Вестника

\renewcommand{\VestnikMonth}{Ноябрь} % Месяц выхода Вестника

\renewcommand{\VestnikData}{01.11.09} % Дата подписания Вестника в печать

\renewcommand{\VestnikUPL}{26,64} % Усл. печ. л.

\renewcommand{\VestnikUIL}{25,73} % Уч.-изд. л.

\renewcommand{\VestnikRegNumber}{479/09} % Регистрационный номер

\renewcommand{\VestnikOrder}{994} % Номер заказа







%

%\begin{document}

%\allowdisplaybreaks[4]

%

%\tableofcontents

%\newpage



\renewcommand{\firstpage}{226}

\renewcommand{\lastpage}{230}







\udk{621.317} \msc{42А10}





\title{Оценка погрешности аппроксимационного метода измерения интегральных характеристик по~отдельным мгновенным значениям

сигналов}% Полный заголовок

{Оценка погрешности аппроксимационного метода \dots}% Заголовок, который идёт в верхний колонтитул



\author{В.\,С.~Мелентьев, Е.\,Е.~Ярославкина, А.\,Н.~Болотнова}{М\,е\,л\,е\,н\,т\,ь\,е\,в~В.\,С., Я\,р\,о\,с\,л\,а\,в\,к\,и\,н\,а~Е.\,Е.,

Б\,о\,л\,о\,т\,н\,о\,в\,а~А.\,Н. }





\institute{\samgtu.}



\email{E-mail}{vs\_mel@mail.ru}



\otherinf{\fullauthor{Владимир Сергеевич Мелентьев} \regalia{д.т.н., доцент} \position{зав. кафедрой}

\dept{каф. информационно"=измерительной техники}

\fullauthor{Екатерина Евгеньевна Ярославкина} \position{аспирант} \dept{каф. информационно"=измерительной техники}

\fullauthor{Анна Николаевна Болотнова} \position{аспирант} \dept{каф. информационно"=измерительной техники}}



\keywords{гармоническая модель, результирующая погрешность, интегральные характеристики}



\workabstract{Рассматриваются различные подходы к оценке результирующей погрешности определения интегральных характеристик

сигналов по отдельным мгновенным значениям, обусловленной отклонением реальных сигналов от используемой гармонической модели.}



\engtitle{Estimation of Error of~Approximation Method of~Measurement of~Integrated Characteristics on~the Separate Instant Values Connected with Transitions of Signals Through the Zero Point}



\engauthor{V.\,S.~Melentyev, E.\,E.~Yaroslavkina, A.\,N.~Bolotnova}{M\,e\,l\,e\,n\,t\,y\,e\,v~V.\,S.,

Y\,a\,r\,o\,s\,l\,a\,v\,k\,i\,n\,a~E.\,E., B\,o\,l\,o\,t\,n\,o\,v\,a~A.\,N.}





\otherenginf{\fullauthor{Vladimir S. Melentyev} \regalia{Dr. Sci. (Techn.)}

\position{Head of Dept.} \dept{Dept. of Information Measuring Systems}

\fullauthor{Ekaterina E. Yaroslavkina} \position{Postgraduate Student}\dept{Dept. of Information Measu\-ring Systems}

\fullauthor{Anna N. Bolotnova} \position{Postgraduate Student}\dept{Dept. of Information Measuring Systems}}



\enginstitute{\engsamgtu.}



\engabstract{Approaches to the estimation of resulting inaccuracy of definition of integrated characteristics of signals on the separate instant values, caused by a deviation of real signals from used harmonious model are studied.}



\engkeywords{harmonious model, resultant an error, integrated characteristics}



\date{14/I/2010}

\revisiondate{15/III/2010}



\maketitle \small



В настоящее время широкое распространение получили методы определения интегральных характеристик гармонических сигналов по их

мгновенным значениям, обеспечивающие время измерения менее периода сигнала. Упрощение реализации данных методов обеспечивает

использование перехода сигналов через ноль. Модели большинства реальных сигналов отличаются от гармонической. Возникает задача

оценки погрешности результата измерения интегральных характеристик при искажении формы сигналов.



При малых углах сдвига фаз между напряжением и током сокращение времени измерения обеспечивает разработанный метод, согласно которому интегральные

характеристики сигналов определяются по двум мгновенным значениям напряжения и одному мгновенному значению тока. Для данного

метода [\citen{mel10:Batishev}] мгновенные значения определяются соотношениями

\[

%\left\{ {\begin{array}{l}

 I_1 =I_m \sin \left( -\varphi  \right), \quad

 U_2 =U_m \sin \varphi , \quad

 U_3 =U_m \sin \left( 2\varphi \right), 

% \\

% \end{array}} \right.

\]

где $I_{1}$ "--- мгновенное значение тока, взятое в момент перехода сигнала напряжения через ноль; $U_{2}$ "--- мгновенное

значение напряжения, взятое в момент перехода сигнала тока через ноль; $U_{3}$ "--- мгновенное значение напряжения, взятое через

промежуток времени, равный интервалу времени между моментами перехода через ноль сигналов напряжения и тока; $U_m$, $I_m$ "---

амплитудные значения сигналов напряжения и тока; $\varphi$ "--- угол сдвига фаз между сигналами напряжения и~тока.



Для гармонических моделей напряжения и тока выражения для определения

интегральных характеристик имеют следующий вид:

\begin{equation}

\label{mel10:eq1}

U_\text{СКЗ} = \frac{2U_2^2}{\sqrt {2\left( 4U_2^2 -U_3^2 \right)}},\quad

%\label{mel10:eq2}

I_\text{СКЗ} =\frac{2|I_1 U_2|}{\sqrt {2\left( 4U_2^2 -U_3^2 \right)}

}

\end{equation}

"--- среднеквадратические значения напряжения и тока;



\begin{equation}

%\label{mel10:eq3}

P=\text{sign}\left( U_2 \right) \text{sign} \left( U_3 \right)\cdot \left| \frac{I_1 U_2^2

U_3}{4U_2^2 - U_3^2}\right|,\quad

Q=-\frac{I_1 U_2^2}{\sqrt {4U_2^2 - U_3^2}}

\label{mel10:eq4}

\end{equation}

"--- активная и реактивная мощности соответственно.



При малых углах сдвига фаз $\varphi$ между сигналами напряжения и тока рассматриваемый

метод обеспечивает малое время измерения $\Delta t_u = 2 \Delta t_\varphi

+\Delta t_{\text{н}} $, где \mbox{$\Delta t_\varphi{=}{\varphi}/{\omega}$}; $\Delta

t_{\text{н}}$ "--- промежуток времени с момента начала измерения до момента перехода

сигнала тока через ноль.



Рассматриваемый метод предназначен для определения интегральных

характеристик в цепях с гармоническими напряжениями и токами. Поэтому

необходимо оценить погрешность метода из-за несоответствия модели реальному

сигналу [\citen{mel10:Melentyev}].



В общем случае при наличии высших гармоник сигналы напряжения и тока имеют

вид

\begin{equation}

\label{mel10:eq5}

u\left( t \right)=\sum \limits_{k=1}^\infty U_{mk} \sin \left( k\omega

t+\psi _{uk} \right),\quad 

%\label{mel10:eq6}

i\left( t \right)=\sum\limits_{k=1}^\infty I_{mk} \sin \left( k\omega

t+\psi_{ik} \right),

\end{equation}

где $U_{mk}$, $I_{mk}$ "--- амплитудные значения $k$-тых гармонических

составляющих напряжения и тока; $\psi_{uk}$, $\psi_{ik}$ "--- начальные фазы гармоник напряжения и тока $k$-того порядка.



Сдвиги переходов через ноль $\Delta \beta_{u}$ и $\Delta \beta_{i}$ сигналов

напряжения и тока относительно их первых гармоник, обусловленные наличием

высших гармоник, можно найти, приравняв к нулю выражения (\ref{mel10:eq5}):

\begin{gather}

\label{mel10:eq7}

\Delta \beta _u =\pm \arcsin \sum\limits_{k=2}^\infty h_{uk} \sin \left(

k\omega t+\psi _{1uk} \right),\quad

%\label{mel10:eq8}

\Delta \beta _i =\pm \arcsin \sum\limits_{k=2}^\infty h_{ik} \sin \left(

k\omega t + \psi_{1ik} \right),

\end{gather}

где $h_{uk}$, $h_{ik}$ "--- коэффициенты гармоник напряжения и тока $k$-того порядка;

$h_{uk} = {U_{mk}}/{U_{m1}}$, $h_{ik} = {I_{mk}}/{I_{m1}}$.



Выражения (\ref{mel10:eq7}) принимают максимальные значения 

$$\Delta \beta_{u_{\max}} = {\pm}\arcsin \sum \limits_{k=2}^\infty h_{uk}, \quad \Delta \beta

_{i_{\max}} = \pm \arcsin \sum \limits_{k=2}^\infty {h_{ik}}$$ при 

$k\omega t + \psi_{1uk} = (2n+1)\pi/2$,  $k \omega t + \psi_{1ik} = (2n+1){\pi}/{2}$,

где  $n\in \mathbb N \cup \{0\}$.



Таким образом, сдвиги переходов через ноль максимальны в том случае, когда

начальные фазы высших гармоник напряжения и тока относительно первых

гармоник $\psi_{1uk}$, $\psi_{1ik}$ кратны $(2n+1){\pi}/{2}$.



Если угол сдвига фаз между сигналами напряжения и тока $\varphi =\varphi_1

+\Delta \beta _0$, где $\Delta \beta_{0} = \Delta \beta_{u}+\Delta \beta_{i}$, то

максимальное его значение имеет место при

\begin{equation}

\label{mel10:eq9}

\Delta \beta _{0\max} =\left| \Delta \beta_{u \max}\right| +

\left| \Delta \beta_{i\max} \right|.

\end{equation}



В общем случае мгновенные значения сигналов напряжения и тока имеют

следующий вид:

\begin{gather}

%\label{mel10:eq10}

\notag

I_{11}=I_{1m} \left\{ \sin \left( -\varphi_1 -\Delta \beta _u \right) + \sum \limits_{k=2}^\infty h_{ik} \sin \Bigl[ k \left( -\varphi _1

+ \Delta \beta_u \right)+\psi_{1ik} \Bigr] \right\};\\

\label{mel10:eq12}

U_{12} =U_{1m} \left\{ \sin \left( \varphi_1 + \Delta \beta_i \right) + \sum\limits_{k=2}^\infty h_{uk} \sin \Bigl[ k\left( \varphi_1

+ \Delta \beta _i \right) + \psi_{1uk} \Bigr] \right\}; \\

%\label{mel10:eq12}

\notag

U_{13} = U_{1m} \left\{ \sin \Bigl[ 2 \left( \varphi_1 {+} \Delta \beta_i \right){+}\Delta \beta _u \Bigr]{+}\sum\limits_{k=2}^\infty h_{uk}

\sin \Bigl[ k \left( 2 \varphi_1 {+} 2 \Delta \beta_i {+} \Delta \beta_u \right) {+} \psi _{1uk} \Bigr] \right\}.

\end{gather}

Подставляя (\ref{mel10:eq12}) в выражения (\ref{mel10:eq1}), (\ref{mel10:eq4}), можно определить интегральные

характеристики сигналов с учётом погрешности из"=за отклонения реальных

сигналов от гармонической модели. Используя полученные выражения, можно

определить относительные погрешности измерения среднеквадратических значений напряжения и тока, а

также приведенные погрешности измерения активной и реактивной мощностей в~соответствии с~выражениями

\begin{gather*}

\delta _{U_1} = \frac{U_\text{СКЗ} - U_\text{СКЗр}}{U_\text{СКЗр}}, \quad \delta _{I_1} = \frac{I_\text{СКЗ} - I_\text{СКЗр}}{I_\text{СКЗр}},\quad

\gamma _{P_1} = \frac{P-P_{\rm p}}{S}, \quad \gamma _{Q_1} =\frac{Q-Q_{\rm p} }{S},

\end{gather*}

$S=U_\text{СКЗ} I_\text{СКЗ}$ "--- полная мощность, а $U_\text{СКЗр}$, $I_\text{СКЗр}$, $P_{\rm p}$, $Q_{\rm p}$ "--- расчётные значения среднеквадратических значений напряжения, тока, активной и реактивной мощностей, соответственно вычисляемые как

\begin{gather*}

U_\text{СКЗр} =\frac{U_{m1}}{\sqrt 2} \sqrt {1+\sum_{k=2}^\infty

{h_{uk}^2}}, \quad I_\text{СКЗр} =\frac{I_{m1} }{\sqrt 2 }\sqrt {1+\sum_{k=2}^\infty

{h_{ik}^2 }},\\

P_{\rm p} =\frac{U_{m1} I_{m1} }{2}\Bigl[ \cos \varphi _1

+\sum\limits_{k=2}^\infty h_{uk} h_{ik} \cos \left( \psi _{uk} -\psi _{ik}

\right) \Bigr],\\

Q_{\rm p} =\frac{U_{m1} I_{m1}}{2}\Bigl[ \cos \varphi _1

+\sum\limits_{k=2}^\infty h_{uk} h_{ik} \sin \left( \psi_{uk} -\psi_{ik}

\right)\Bigr].

\end{gather*}





Проведённое моделирование зависимости погрешностей определения интегральных

характеристик сигналов $\delta_{U_1}$, $\delta_{I_1}$, $\gamma_{P_1}$ и $\gamma_{Q_1}$ от угла сдвига

фаз между первыми гармониками напряжения и тока $\varphi_{1}$ при различном

гармоническом составе сигналов показало, что при искажении формы сигналов

появляется существенная погрешность, зависящая от угла сдвига фаз между

первыми гармониками напряжения и тока и спектрального состава сигналов.



Оценка данного вида погрешности определения интегральных характеристик по

расчетным значениям дает наиболее точный результат. Однако для этого

необходимо знать вид реальных сигналов.



Более просто можно оценить погрешность из-за отклонения реальных сигналов от

гармонической модели с помощью метода [\citen{mel10:Melentyev}], основанного на определении

погрешности измерения интегральной характеристики как функции, аргументы

которой заданы приближенно с погрешностью, соответствующей максимальному

отклонению модели от реального сигнала.



При этом абсолютные погрешности определения среднеквадратических значений напряжения и тока, активной и реактивной мощностей

рассчитываются из выражений

\begin{equation}

\label{mel10:eq13}

\begin{array}{l}

\Delta U_{СКЗ} =\left| ( U_{СКЗ} )^\prime_{U_2}\right| \cdot  \Delta U_2 + \left| ( U_{СКЗ} )^\prime_{U_3}\right| \cdot \Delta U_3, \vspace{1mm}\\

%\label{mel10:eq14}

\Delta I_{СКЗ} =\left| ( I_{СКЗ}  )^\prime _{I_1}\right| \cdot \Delta I_1 +\left| ( I_{СКЗ} )^\prime _{U_2}\right| \cdot  \Delta U_2 +\left| ( I_{СКЗ} )^\prime_{U_3} \right| \cdot \Delta U_3,\vspace{1mm}\\

%\label{mel10:eq15}

\Delta P=\left| ( P )^\prime _{I_1} \right| \cdot \Delta I_1 +\left| ( P )^\prime _{U_2} \right| \cdot \Delta U_2 + 

\left| ( P)^\prime_{U_3} \right| \cdot \Delta U_3,\vspace{1mm}\\

%\label{mel10:eq16}

\Delta Q=\left|( Q )^\prime _{I_1}  \right| \cdot \Delta I_1 +\left|( Q)^\prime _{U_2} \right| \cdot \Delta U_2 +\left| ( Q)^\prime_{U_3} \right| \cdot \Delta U_3.

\end{array}

\end{equation}

После вычисления производных выражения (\ref{mel10:eq13}) примут вид

\begin{equation}

\label{mel10:eq17}

\begin{array}{l}

\Delta U_\text{СКЗ} = \frac{\sqrt 2}{\left| \sin^3 \varphi_1 \right|} \left(

\left| \sin^2 \varphi_1 - 0,5 \right|\Delta U_2 + \left| 0{,}25 \cos \varphi

_1  \right| \Delta U_3 \right), \vspace{1mm} \\

%\label{mel10:eq18}

\Delta I_\text{СКЗ} =\frac{\sqrt 2}{2\left| \sin \varphi_1 \right|} \left(

\Delta I_1 + \frac{I_{1m}}{U_{1m}} \ctg^2 \varphi_1 \Delta U_2

+ \frac{I_{1m}}{U_{1m}} \left| \frac{\ctg \varphi_1

}{2 \sin \varphi_1} \right| \Delta U_3 \right), \vspace{1mm}\\

%\label{mel10:eq19}

\Delta P= \frac{1}{2} U_{1m} \left| \ctg \varphi_1 \right| \Delta I_1

+I_{1m} \left| \ctg^3 \varphi_1 \right|\Delta U_2 +I_{1m}

\frac{1+\cos^2 \varphi_1} {4 \left|

\sin ^3 \varphi_1 \right|} \Delta U_3, \vspace{1mm}\\

%\label{mel10:eq20}

\Delta Q=\frac{1}{2}\left( U_{1m} \Delta I_1 +I_{1m} \left|

1-\ctg ^2 \varphi_1 \right| \Delta U_2 +I_{1m} \left|

\frac{\ctg \varphi_1}{2 \sin \varphi _1} \right| \Delta U_3 \right).

\end{array}

\end{equation}



С учётом (\ref{mel10:eq17}) относительные погрешности определения среднеквадратических значений напряжения и

тока и приведённые погрешности определения активной и реактивной мощностей определяются как

\begin{equation}

\label{mel10:eq21}

\begin{array}{l}

\delta _{U_\text{СКЗ\/2}} = \frac{2}{U_{1m} \left| \sin^3 \varphi_1

\right|} \left( \left| \sin^2 \varphi_1 - 0,5 \right| \Delta U_2 + \left|

0,25 \cos \varphi_1 \right| \Delta U_3 \right),  \vspace{1mm}\\

%\label{mel10:eq22}

\delta _{I_\text{СКЗ\/2}} =\frac{1}{\left| \sin \varphi_1 \right|} \left(

\frac {\Delta I_1}{I_{1m}} + \frac{\Delta U_2}{U_{1m}}\ctg^2 \varphi

_1 + \left| \frac{\ctg \varphi_1 }{2 \sin \varphi_1} \right| \frac{\Delta U_3}{U_{1m}} \right), \vspace{1mm}\\

%\label{mel10:eq23}

\gamma_{P_2} = \frac{1}{I_{1m}} \left| \ctg \varphi_1 \right| \Delta

I_1 + \frac{\Delta U_2}{U_{1m}} \left| \ctg^3 \varphi_1

\right|+\frac{\Delta U_3}{U_{1m}} \frac{1 + \cos^2 \varphi

_1}{4 \left| \sin^3 \varphi_1  \right|}, \vspace{1mm}\\

%\label{mel10:eq24}

\gamma _{Q_2} =\frac{\Delta I_1}{I_{1m}} + \left| 1-\ctg^2 \varphi_1

\right| \frac{\Delta U_2}{U_{1m}} + \left| \frac{\ctg \varphi_1

}{2\sin \varphi_1} \right| \frac{\Delta U_3}{U_{1m}}.

\end{array}

\end{equation}



Абсолютные погрешности аргументов в моменты времени $t_{1}$, $t_{2}$ и

$t_{3}$ определяются выражениями

\begin{equation}

\label{mel10:eq25}

\begin{array}{l}

\Delta I_1 = \sum\limits_{k=1}^\infty I_{km} \sin \left[ k \left( -\varphi

_1 -\Delta \beta _u \right)+ \psi_{1ik} \right] -I_{1m} \sin \left(

-\varphi_1 \right), \vspace{1mm}\\

%\label{mel10:eq26}

\Delta U_2 =\sum \limits_{k=1}^\infty U_{km} \sin \left[ k\left( \varphi

_1 + \Delta \beta _i \right)+\psi_{1uk} \right] - U_{1m} \sin \varphi _1,\vspace{1mm}\\

%\label{mel10:eq27}

\Delta U_3 =\sum \limits_{k=1}^\infty U_{km} \sin \left[ k\left( 2\varphi

_1 +2\Delta \beta _i +\Delta \beta _u  \right)+\psi _{1uk} \right]

-U_{1m} \sin \left( 2\varphi _1 \right).

\end{array}

\end{equation}



Если считать, что абсолютные погрешности аргументов соответствуют наибольшим

отклонениям в моменты времени $t_{1}$, $t_{2}$ и $t_{3}$, то (\ref{mel10:eq25}) с уч\"етом (\ref{mel10:eq7})--(\ref{mel10:eq12}) приводятся к виду

\begin{equation}

\label{mel10:eq28}

\begin{array}{l}

\sup \left( \Delta I_1 \right)=I_{1m} \left| \sin \left( -\varphi_1

- \arcsin \sum \limits_{k=2}^\infty h_{uk} \right)+ \sum \limits_{k=2}^\infty h_{ik} - \sin \left( -\varphi _1

\right) \right|, \vspace{1mm} \\

%\label{mel10:eq29}

\sup \left( \Delta U_2 \right)=U_{1m} \left| \sin \left( \varphi_1

+ \arcsin \sum\limits_{k=2}^\infty h_{ik} \right)+\sum\limits_{k=2}^\infty h_{uk} - \sin \varphi_1 \right|, \vspace{1mm} \\

%\end{gather}

%\begin{multline}

%\label{mel10:eq30}

\sup \left( \Delta U_3  \right){=}U_{1m} 

%\times\\

%\times 

\left| \sin \left( 2\varphi _1

{+}2 \arcsin \sum \limits_{k=2}^\infty h_{ik}{+}\arcsin

\sum \limits_{k=2}^\infty h_{uk} \right){+}\sum\limits_{k=2}^\infty

h_{uk}{-}\sin \left( 2\varphi_1 \right)\right|. \hspace{-5mm}

\end{array}

\end{equation}





По соотношениям (\ref{mel10:eq21}) с уч\"етом (\ref{mel10:eq28}) было проведено

моделирование зависимости погрешностей определения интегральных характеристик сигналов $\delta_{U_2}$, $\delta_{I_2}$,

$\gamma_{P_2}$ и $\gamma_{Q_2}$ от угла сдвига фаз между первыми гармониками напряжения и тока $\varphi_{1}$ при различном

гармоническом составе сигналов. Анализ показал, что при искажении формы сигналов появляется существенная погрешность, зависящая

от угла сдвига фаз между первыми гармониками напряжения и тока. Величина погрешности возрастает с ростом номера и коэффициента

гармоники и может достигать $10\%$ при $h_{uk}\approx h_{ik}\approx 0{,}01$. Второй метод оценки погрешности результата измерения интегральной

характеристики как функции, аргументы которой заданы приближенно с погрешностью, соответствующей наибольшему отклонению модели от

реального сигнала в соответствующих точках, дает завышенные значения погрешности, особенно для малых углов сдвига фазы

$\varphi_{1}$.



Полученные результаты позволяют принимать решение о возможности использования разработанного аппроксимационного метода для

определения интегральных характеристик (среднеквадратические значения напряжения и тока, активной и реактивной мощности) искаженных сигналов в зависимости

от предъявляемых требований по точности и быстродействию.





\begin{thebibliography}{99}



\RBibitem{mel10:Batishev}

\by Батищев~В.\,И., Мелентьев~В.\,С.

\book Аппроксимационные методы и~системы промышленных измерений, контроля, испытаний, диагностики

\yr 2007

\publaddr М.

\publ Машиностроение-1

\finalinfo 393~c



\RBibitem{mel10:Melentyev}

\by Мелентьев~В.\,С.

\paper Методы оценки влияния погрешностей, обусловленных не соответствием модели виду реального сигнала, на погрешность результата

измерения

\jour Вестн. Сам. гос. техн. ун-та. Сер. Техн. науки

\yr 2006

\issue 41

\pages 89--96



\end{thebibliography}



\makeengtitle





%\end{document}

